\chapter{The Stone Piano Hypothesis: An Archaeological Case Study}
\label{app:stone piano}

\epigraph{Whatever they were doing, they intended it.}{}

The Stone Piano Hypothesis proposes that the Bruniquel Cave stalagmite structures,
dated to approximately 176,500 years before present 	tep{jaubert2016} and attributed to \emph{Homo
neanderthalensis}, functioned as lithophones: percussion instruments fabricated
from resonant calcite speleothems. This appendix situates the hypothesis within
the Yarncrawler framework.

\section{Admissibility Filtering in Speleothem Selection}

The fundamental frequency of a free-free calcite rod of length $L$ and radius $r$:
\[
  f 1 = \frac{\lambda 1^2\, r}{8\pi L^2}\sqrt{\frac{E}{\rho}},
  \qquad \lambda 1 \approx 4.730,
\]
with $E \approx 75\,\text{GPa}$ and $\rho \approx 2600\,\text{kg\,m}^{-3}$.
Fracturing stalagmites from their bases converts clamped-free resonators
(inharmonic series $1 : 6.27 : 17.55$) into free-free resonators with more harmonic
profiles ($1 : 2.76 : 5.40$). This is an admissibility filtering operation.

\section{Sympathetic Resonance as the Corrigibility Mechanism}

Acoustic beating---the periodic amplitude modulation produced by two nearly
coincident frequencies---disappears when frequencies are matched. Its disappearance
is an unambiguous perceptual event providing real-time corrigibility feedback. The
builders could tune fragment lengths by progressive shortening, without abstract
acoustic knowledge, using only the disappearance of beating as a signal. This is the
\KES loop operating in stone at 176,500 BP.

\section{The Ring as Stigmergic Structure}

A partially assembled ring of resonant elements constitutes an acoustic admissibility
filter. New candidate fragments that excite strong sympathetic vibrations in the
existing ring are harmonically admissible; those that do not are refused. The
structure under construction participates in selecting the components from which it
is built: Axiom~3 (Autocatalytic Fix Design) instantiated in deep prehistory.

\section{Estimated Chamber Reverberation}

By Sabine's formula with $V \approx 73\,\text{m}^3$ and $\alpha \approx 0.03$:
\[
  T {60} \approx \frac{0.161 \times 73}{0.03 \times 130} \approx 3.0\,\text{s}.
\]
The cave amplifies and extends each repair operation, lowering the threshold for
subsequent ones. It is a Yarncrawler environment.
